\uuid{vGQU}
\exo7id{3837}
\titre{exo7 3837}
\auteur{quercia}
\organisation{exo7}
\datecreate{2010-03-11}
\isIndication{false}
\isCorrection{false}
\chapitre{Espace euclidien, espace normé}
\sousChapitre{Espaces vectoriels hermitiens}
\cpge{mpsi-mp2i.espaces-prehilbertiens.a, pcsi.espaces-prehilbertiens.a, ptsi.geometrie-plan.b, mpsi-mp2i.derivabilite.d, pcsi.derivabilite.d, ptsi.derivabilite.d, mpsi-mp2i.espaces-vectoriels.a, pcsi.espaces-vectoriels.a, ptsi.geometrie-espace.c, bcpst1.geometrie.f, bcpst1.derivation.c, psi.espaces-prehilbertiens.a, pt.structure-prehilbertienne.a, mp-mpi.topologie-espaces-normes.a, pc.espaces-vectoriels-normes.a, psi.espaces-vectoriels-normes.a, bcpst2.produit-scalaire.a}
\module{Algèbre}
\niveau{L2}
\difficulte{}

\contenu{
\question{
Soit $D$ le disque unité fermé de $\R^2$.
On considère l'espace $E$ des fonctions $f : D \to \R$ de classe $\mathcal{C}^1$
nulles sur le bord, $C$, de $D$.

Pour $f,g \in E$, on pose $(f\mid g) = \iint_D \Bigl(
\frac{\partial f}{\partial x}
\frac{\partial g}{\partial x} +
\frac{\partial f}{\partial y}
\frac{\partial g}{\partial y}
\Bigr) d x d y$.
Montrer que c'est un produit scalaire.
}
}
